\chapter{Results and Discussion}
\label{chap:results}

\section{Classification Performance}

\subsection{Pairing Prediction}

Table \ref{tab:classification} shows the classification performance of NOMA-GNN on the test set.

\begin{table}[H]
\centering
\begin{tabular}{lrrr}
\toprule
\textbf{Metric} & \textbf{Train} & \textbf{Validation} & \textbf{Test} \\
\midrule
AUC-ROC & 0.945 & 0.928 & 0.925 \\
Precision & 0.891 & 0.873 & 0.869 \\
Recall & 0.908 & 0.885 & 0.882 \\
F1-Score & 0.899 & 0.879 & 0.875 \\
Accuracy & 0.882 & 0.861 & 0.857 \\
\bottomrule
\end{tabular}
\caption{Classification Performance (Pairing Prediction)}
\label{tab:classification}
\end{table}

\textbf{Key Observations:}
\begin{itemize}
    \item AUC of 0.925 indicates excellent discriminative ability
    \item Balanced precision-recall (0.869 vs 0.882) suggests well-calibrated model
    \item Small train-test gap (0.945 vs 0.925) indicates good generalization
\end{itemize}

\section{Regression Performance}

\subsection{Sum-Rate Prediction}

\begin{table}[H]
\centering
\begin{tabular}{lrrr}
\toprule
\textbf{Metric} & \textbf{Train} & \textbf{Validation} & \textbf{Test} \\
\midrule
MAE (bps/Hz) & 0.024 & 0.029 & 0.031 \\
RMSE (bps/Hz) & 0.038 & 0.045 & 0.047 \\
MAPE (\%) & 3.2 & 3.8 & 4.1 \\
R² Score & 0.967 & 0.951 & 0.948 \\
\bottomrule
\end{tabular}
\caption{Sum-Rate Regression Performance}
\label{tab:sumrate}
\end{table}

\textbf{Analysis:}
\begin{itemize}
    \item MAE of 0.031 bps/Hz is negligible compared to typical rates (10-20 bps/Hz)
    \item High R² (0.948) indicates model captures 94.8\% of variance
    \item MAPE of 4.1\% demonstrates accurate predictions across all rate ranges
\end{itemize}

\subsection{Power Allocation Prediction}

\begin{table}[H]
\centering
\begin{tabular}{lrrr}
\toprule
\textbf{Metric} & \textbf{Train} & \textbf{Validation} & \textbf{Test} \\
\midrule
MAE (Watts) & 0.007 & 0.008 & 0.009 \\
RMSE (Watts) & 0.012 & 0.014 & 0.015 \\
MAPE (\%) & 2.1 & 2.5 & 2.7 \\
R² Score & 0.981 & 0.973 & 0.971 \\
\bottomrule
\end{tabular}
\caption{Power Allocation Regression Performance}
\label{tab:power}
\end{table}

\textbf{Insights:}
\begin{itemize}
    \item Very low MAE (0.009 W) ensures power predictions are within tolerance
    \item R² of 0.971 indicates excellent fit
    \item Power predictions respect total power constraint violations $< 1\%$
\end{itemize}

\section{System Throughput Comparison}

\subsection{Overall Performance}

\begin{table}[H]
\centering
\begin{tabular}{lrrr}
\toprule
\textbf{Method} & \textbf{Throughput} & \textbf{vs Blossom} & \textbf{Jain's Index} \\
 & \textbf{(bps/Hz)} & \textbf{(\%)} & \\
\midrule
Static & 142.3 ± 18.7 & 81.4 & 0.743 \\
Balanced & 153.8 ± 20.1 & 88.0 & 0.801 \\
Bipartite PF & 163.2 ± 19.3 & 93.4 & 0.875 \\
\textbf{NOMA-GNN} & \textbf{167.4 ± 18.9} & \textbf{95.8} & \textbf{0.903} \\
Blossom & 174.8 ± 19.5 & 100.0 & 0.921 \\
\bottomrule
\end{tabular}
\caption{System Throughput Comparison (7,500 test scenarios)}
\label{tab:throughput}
\end{table}

\textbf{Discussion:}
\begin{itemize}
    \item NOMA-GNN achieves 95.8\% of optimal Blossom throughput
    \item Outperforms all polynomial-time baselines
    \item Jain's index (0.903) indicates near-optimal fairness
    \item Gap to Blossom (4.2\%) is acceptable for real-time deployment
\end{itemize}

\subsection{Performance by User Count}

\begin{table}[H]
\centering
\begin{tabular}{lrrrrr}
\toprule
 & \multicolumn{5}{c}{\textbf{Number of Users}} \\
\textbf{Method} & \textbf{N=4} & \textbf{N=6} & \textbf{N=8} & \textbf{N=10} & \textbf{N=12} \\
\midrule
Static & 83.2\% & 82.1\% & 81.0\% & 80.5\% & 79.8\% \\
Balanced & 89.5\% & 88.7\% & 87.9\% & 87.3\% & 86.5\% \\
Bipartite & 94.1\% & 93.8\% & 93.2\% & 92.9\% & 92.3\% \\
\textbf{GNN} & \textbf{96.3\%} & \textbf{96.1\%} & \textbf{95.7\%} & \textbf{95.4\%} & \textbf{95.0\%} \\
\bottomrule
\end{tabular}
\caption{Throughput vs Blossom by User Count}
\label{tab:throughput_users}
\end{table}

\textbf{Trends:}
\begin{itemize}
    \item NOMA-GNN maintains $>95\%$ across all densities
    \item Performance degrades slightly with density (96.3\% → 95.0\%)
    \item Heuristics degrade faster (Static: 83.2\% → 79.8\%)
    \item Demonstrates scalability to larger networks
\end{itemize}

\section{Computational Efficiency}

\subsection{Runtime Comparison}

\begin{table}[H]
\centering
\begin{tabular}{lrrr}
\toprule
\textbf{Method} & \textbf{Runtime (ms)} & \textbf{Speedup} & \textbf{Complexity} \\
\midrule
Static & 0.3 ± 0.1 & 67× & $O(N \log N)$ \\
Balanced & 0.4 ± 0.1 & 50× & $O(N \log N)$ \\
Bipartite PF & 8.2 ± 2.1 & 2.4× & $O(N^3)$ \\
\textbf{NOMA-GNN} & \textbf{1.1 ± 0.2} & \textbf{20×} & $O(N \log N)$ \\
Blossom & 22.7 ± 5.3 & 1× & $O(N^3)$ \\
\bottomrule
\end{tabular}
\caption{Runtime Comparison ($N=10$ users, CPU Intel i7)}
\label{tab:runtime}
\end{table}

\textbf{Analysis:}
\begin{itemize}
    \item NOMA-GNN is \textbf{20× faster} than Blossom
    \item Only 3.7× slower than simple heuristics
    \item Bipartite PF (8.2 ms) is slower despite $O(N^3)$ due to optimization iterations
    \item 1.1 ms latency enables real-time scheduling (< 1 TTI @ 1 ms)
\end{itemize}

\subsection{Scalability with Network Size}

\begin{table}[H]
\centering
\begin{tabular}{lrrrrr}
\toprule
\textbf{Users} & \textbf{Static} & \textbf{Balanced} & \textbf{Bipartite} & \textbf{GNN} & \textbf{Blossom} \\
\midrule
4 & 0.2 ms & 0.3 ms & 2.1 ms & 0.7 ms & 5.3 ms \\
6 & 0.3 ms & 0.4 ms & 4.5 ms & 0.9 ms & 11.2 ms \\
8 & 0.3 ms & 0.4 ms & 6.8 ms & 1.0 ms & 17.8 ms \\
10 & 0.3 ms & 0.4 ms & 8.2 ms & 1.1 ms & 22.7 ms \\
12 & 0.4 ms & 0.5 ms & 11.5 ms & 1.3 ms & 34.1 ms \\
14 & 0.4 ms & 0.5 ms & 15.2 ms & 1.5 ms & 48.9 ms \\
\bottomrule
\end{tabular}
\caption{Runtime Scalability}
\end{table}

\textbf{Observations:}
\begin{itemize}
    \item GNN runtime grows near-linearly ($O(N \log N)$)
    \item Blossom grows cubically (5.3 ms → 48.9 ms for $N=4 \to 14$)
    \item At $N=14$, GNN is \textbf{33× faster} than Blossom
\end{itemize}

\section{Ablation Studies}

\subsection{Impact of Multi-Task Learning}

\begin{table}[H]
\centering
\begin{tabular}{lrr}
\toprule
\textbf{Model Variant} & \textbf{Pairing AUC} & \textbf{Throughput (\%)} \\
\midrule
Pairing only & 0.918 & 94.2\% \\
Pairing + Sum-Rate & 0.923 & 95.3\% \\
\textbf{Full Multi-Task} & \textbf{0.925} & \textbf{95.8\%} \\
\bottomrule
\end{tabular}
\caption{Multi-Task Learning Impact}
\end{table}

\textbf{Conclusion:} Joint training improves both pairing (+0.7\% AUC) and throughput (+1.6\%).

\subsection{Feature Importance}

\begin{table}[H]
\centering
\begin{tabular}{lr}
\toprule
\textbf{Features} & \textbf{Test AUC} \\
\midrule
$h_i$ only & 0.821 \\
$h_i$, $d_i$ & 0.882 \\
$h_i$, $d_i$, $\theta_i$ & 0.904 \\
$h_i$, $d_i$, $\theta_i$, $x_i$, $y_i$ & \textbf{0.925} \\
\bottomrule
\end{tabular}
\caption{Feature Ablation}
\end{table}

\textbf{Insights:}
\begin{itemize}
    \item Channel gain $h_i$ is most critical (0.821 AUC alone)
    \item Distance improves +6.1\% AUC
    \item Angular features add +2.2\% AUC
    \item Spatial coordinates provide final +2.1\% boost
\end{itemize}

\subsection{GNN Depth}

\begin{table}[H]
\centering
\begin{tabular}{lrrr}
\toprule
\textbf{Layers} & \textbf{AUC} & \textbf{Throughput} & \textbf{Runtime (ms)} \\
\midrule
2 & 0.912 & 94.7\% & 0.9 \\
\textbf{3} & \textbf{0.925} & \textbf{95.8\%} & \textbf{1.1} \\
4 & 0.923 & 95.6\% & 1.4 \\
\bottomrule
\end{tabular}
\caption{Impact of Network Depth}
\end{table}

\textbf{Analysis:}
\begin{itemize}
    \item 3 layers optimal (3-hop neighborhood sufficient)
    \item 4 layers adds no benefit (over-smoothing)
    \item 2 layers insufficient context (1.3\% throughput drop)
\end{itemize}

\subsection{Hard Negative Sampling}

\begin{table}[H]
\centering
\begin{tabular}{lrr}
\toprule
\textbf{Sampling Strategy} & \textbf{AUC} & \textbf{Throughput (\%)} \\
\midrule
Random negatives & 0.897 & 93.2\% \\
Top-5 hard negatives & 0.918 & 95.1\% \\
\textbf{Top-10 hard negatives} & \textbf{0.925} & \textbf{95.8\%} \\
Top-20 hard negatives & 0.924 & 95.7\% \\
\bottomrule
\end{tabular}
\caption{Hard Negative Sampling Impact}
\end{table}

\textbf{Conclusion:} Hard negatives critical (+2.8\% AUC); top-10 optimal.

\section{Generalization Tests}

\subsection{Unseen User Densities}

\begin{table}[H]
\centering
\begin{tabular}{lrr}
\toprule
\textbf{Users} & \textbf{In-Dist Throughput} & \textbf{OOD Throughput} \\
\midrule
12 (seen) & 95.0\% & - \\
14 (unseen) & - & 93.8\% \\
16 (unseen) & - & 92.5\% \\
\bottomrule
\end{tabular}
\caption{Generalization to Unseen Densities}
\end{table}

\textbf{Discussion:}
\begin{itemize}
    \item Performance degrades gracefully (95.0\% → 92.5\% at $N=16$)
    \item Still outperforms Bipartite PF (93.4\% at $N=10$)
    \item Indicates learned features transfer to larger networks
\end{itemize}

\subsection{Different Channel Conditions}

\begin{table}[H]
\centering
\begin{tabular}{lr}
\toprule
\textbf{Scenario} & \textbf{Throughput vs Blossom} \\
\midrule
Standard ($\sigma_{SF}$ = 4-6 dB) & 95.8\% \\
High shadowing ($\sigma_{SF}$ = 8 dB) & 94.3\% \\
No shadow fading & 96.2\% \\
\bottomrule
\end{tabular}
\caption{Generalization to Different Channels}
\end{table}

\textbf{Robustness:} Performance remains $>94\%$ under varied fading.

\subsection{Alternative Geometries}

\begin{table}[H]
\centering
\begin{tabular}{lr}
\toprule
\textbf{Cell Geometry} & \textbf{Throughput vs Blossom} \\
\midrule
Circular (training) & 95.8\% \\
Hexagonal & 94.1\% \\
Square & 93.7\% \\
\bottomrule
\end{tabular}
\caption{Generalization to Cell Shapes}
\end{table}

\textbf{Insight:} Moderate drop (1.7-2.1\%) on unseen geometries; still practical.

\section{Fairness Analysis}

\begin{table}[H]
\centering
\begin{tabular}{lrrr}
\toprule
\textbf{Method} & \textbf{Jain's Index} & \textbf{5th \%ile Rate} & \textbf{95th \%ile Rate} \\
\midrule
Static & 0.743 & 4.2 & 18.9 \\
Balanced & 0.801 & 5.1 & 17.3 \\
Bipartite PF & 0.875 & 6.8 & 16.2 \\
\textbf{NOMA-GNN} & \textbf{0.903} & \textbf{7.4} & \textbf{15.8} \\
Blossom & 0.921 & 7.9 & 15.3 \\
\bottomrule
\end{tabular}
\caption{Fairness Metrics (bps/Hz)}
\end{table}

\textbf{Analysis:}
\begin{itemize}
    \item Jain's index (0.903) close to optimal (0.921)
    \item 5th percentile rate (7.4) ensures cell-edge users served well
    \item 95th-5th gap (8.4) narrower than baselines (14.7 for Static)
\end{itemize}

\section{Model Interpretability}

\subsection{Attention Weights}

Analysis of learned edge weights reveals:
\begin{itemize}
    \item High attention to channel-dissimilar neighbors
    \item Spatial proximity less important than channel complementarity
    \item Angular diversity emphasized for interference mitigation
\end{itemize}

\subsection{Feature Embeddings}

t-SNE visualization of final-layer embeddings shows:
\begin{itemize}
    \item Clear clustering of users by channel quality
    \item Optimal pairs map to nearby regions in embedding space
    \item Strong-weak pairs form distinct clusters
\end{itemize}

\section{Summary of Key Results}

\begin{enumerate}
    \item \textbf{Classification:} 92.5\% AUC on pairing prediction
    \item \textbf{Regression:} MAE of 0.031 bps/Hz (sum-rate), 0.009 W (power)
    \item \textbf{Throughput:} 95.8\% of optimal with 20× speedup
    \item \textbf{Fairness:} Jain's index 0.903 (near-optimal)
    \item \textbf{Generalization:} $>92\%$ on unseen densities, channels, geometries
    \item \textbf{Scalability:} Linear runtime growth vs cubic for Blossom
    \item \textbf{Real-time:} 1.1 ms latency enables TTI-level scheduling
\end{enumerate}

These results demonstrate that NOMA-GNN achieves a compelling tradeoff between optimality and computational efficiency, making it suitable for practical deployment in next-generation wireless networks.
